\documentclass[10pt,a4paper]{amsart}
\usepackage[margin=19mm]{geometry}
\usepackage{amsmath,amssymb,mathtools}
\usepackage[scr=rsfs,cal=euler]{mathalfa}
\usepackage{xfrac}
\usepackage[unicode,hidelinks,bookmarks=false]{hyperref}
\newcommand{\ie}{i.e.\ }
\newcommand{\cf}{cf.\ }
\newcommand{\resp}{respectively\ }
\newcommand{\Subsection}{\subsection}
\providecommand{\nobreakdash}{\nobreak\hbox{-}}
\setlength{\emergencystretch}{2em}
\multlinegap0pt
\usepackage{bm}
\usepackage{bbm}
\usepackage{dsfont}
\usepackage{xspace}
\usepackage{cancel}
\usepackage{mathtools}
\usepackage{amsthm}
\makeatletter
\@ifundefined{theorem}{\newtheorem{theorem}{Theorem}}{}
\@ifundefined{lemma}{\newtheorem{lemma}[theorem]{Lemma}}{}
\makeatother
\makeatletter
\DeclareMathOperator{\Span}{span}
\expandafter\ifx\csname mythm\endcsname\relax
\newenvironment{mythm}[1]
  {\renewcommand\theinnercustomthm{#1}\innercustomthm}
  {\endinnercustomthm}
\fi
\makeatother
\title[Graham's rearrangement for a class of semidirect products]{Graham's rearrangement for a class of semidirect products}
\author{Simone Costa, Stefano Della Fiore, Eva R. Engel}
\date{Source: arXiv:2503.18101v2 (CC BY 4.0)}
\begin{document}
\maketitle

\noindent\emph{Source and licence.} Graham's rearrangement for a class of semidirect products. Version arXiv:2503.18101v2 (CC BY 4.0), distributed under the Creative Commons Attribution 4.0 International licence, as recorded in the arXiv licence metadata for this exact version and shown on the abstract page https://arxiv.org/abs/2503.18101v2. Full rights notice: licenses/arxiv-2503.18101-CC-BY-4.0.txt.\par\smallskip

\noindent\textit{Editorial scope.} Counted unit: the Lemma whose source label is \texttt{lem:small-dim-recitification} in the article (source0.tex lines 204--207, source label \texttt{lem:small-dim-recitification}), delivered together with the article's own complete original proof of it (source0.tex lines 209--246, one \texttt{proof} environment of 38 lines). No mathematics is written, repaired or re-derived: statement and proof are the article's own text line for line. Required context retained uncounted: lem:dissociated-set-generator (lines 170--173); lem:dimension-cardinality2 (lines 184--187). Every definition, display and label of those lines is retained unchanged. Omitted from this package and not redistributed: 1--169, 174--183, 188--203, 208--208, 247--703. Nothing inside a retained excerpt is omitted or altered: every retained body is delivered exactly as the article's lines are written, so no omission receipt of any kind accompanies this package. Citations: the retained excerpts carry no citation command at all, so no bibliography entry of the article is transcribed and no reference is redistributed. Reference adaptation: none was needed and none was made; every reference inside the delivered unit resolves inside it, so no unresolved reference or citation is produced. Printed slips kept literal: no printed word, symbol, hypothesis or number of the counted statement or of its proof was altered. Layout and class: the article's own class is \texttt{amsart} with packages including geometry, amssymb, amsthm, amsmath, mathrsfs, amsbsy, xy, bm, hyperref, tikz, array, float, enumerate, xcolor; the wrapper is the lane's pinned \texttt{amsart} editorial wrapper at 10pt on A4, which restates the theorem-like environments the retained text needs and adds the article's own macro definitions that the retained text needs (\texttt{\textbackslash Span}), with extra wrapper packages (bm, bbm, dsfont, xspace, cancel, mathtools). The delivered numbering is the unit's own. Assets: no retained span contains a figure environment, a graphics-inclusion command, a PDF-page inclusion, a table or a TikZ picture; no image file is copied into this package, the package declares no asset dependency and no image file is redistributed with it. Licence: version arXiv:2503.18101v2 of this article is distributed under the Creative Commons Attribution 4.0 International licence, as recorded in the arXiv licence metadata for this exact version and shown on the abstract page https://arxiv.org/abs/2503.18101v2. A full rights notice is delivered at licenses/arxiv-2503.18101-CC-BY-4.0.txt. Characters: every retained excerpt is pure ASCII, so the delivered engine, XeLaTeX with its default Latin Modern text font, reports no missing character.
\begin{lemma}\label{lem:dissociated-set-generator}
Let $B \subseteq G$ be a finite subset of a not necessarily abelian group.  If $D=\{d_1, \ldots, d_r\}$ is a maximal dissociated subset of $B$, then
$$B \subseteq \Span(D) \vcentcolon= \left\{d_{\sigma(1)}^{\epsilon_{\sigma(1)}} \cdot \cdots \cdot d_{\sigma(r)}^ {\epsilon_{\sigma(r)}}: \sigma \in Sym(r), \epsilon_{\sigma(i)}\in\{-1,0,1\}\right\}.$$
\end{lemma}

\begin{lemma}\label{lem:dimension-cardinality2}
Let $G=\mathbb{Z}_{p} \rtimes_{\varphi} H$ where $H$ is abelian and $\varphi$ satisfies property $(\star)$. Let $B$ be a finite subset of $G$. Then
$|B| \leq 5^{\dim(B)}$.
\end{lemma}

\begin{lemma}\label{lem:small-dim-recitification}

Let $G=\mathbb{Z}_{p} \rtimes_{\varphi} H$ where $H$ is abelian and $\varphi$ satisfies property $(\star)$. If $B \subseteq G$ is a nonempty subset of dimension $\dim(B)< R$, then there is some $\phi \in Aut(G)$ such that the image $\pi_1(\phi(B))$ (where $\pi_1$ is the projection over the $\mathbb{Z}_p$ component) is contained in the interval $(-\frac{p}{100|B|},\frac{p}{100|B|})$.
\end{lemma}

\begin{proof}
Given $\lambda\in \mathbb{Z}_p$, we consider here the family of maps $\phi_{\lambda}: G\rightarrow G$
defined as
$$ \phi_{\lambda}(x,a)=(\lambda x,a).$$
Note that, if $\lambda\not=0$, then $\phi_{\lambda}\in Aut(G)$.

Let us now consider a maximal dissociated subset $D=\{(x_1,a_1),\dots,(x_r,a_r)\}$ of $B$, where $r=|D|= \dim(B)$. Lemma~\ref{lem:dissociated-set-generator} tells us that $B\subseteq \Span(D)$.

Consider the set
$$\{\lambda(x_1,\dots,x_r): \lambda \in \mathbb{Z}_p\} \subseteq (\mathbb{Z}_p)^{r}.$$
Due to the pigeonhole principle, there exists distinct $\lambda_1, \lambda_2 \in \mathbb{Z}_p$ such that $\|\lambda_1 x_i -\lambda_2 x_i\| \leq p^{1-1/r}$ for all $i\in [1,r]$.  Set $\lambda\vcentcolon=\lambda_1-\lambda_2$, we have that $\phi_\lambda\in Aut(G)$ and $\pi_1(\phi_\lambda(D))\subseteq  [-p^{1-1/\dim(B)},p^{1-1/\dim(B)}]$.
Since $B\subseteq \Span(D)$, we have
$$\pi_1(\phi(B)) \subseteq [-\dim(B)p^{1-1/\dim(B)},\dim(B)p^{1-1/\dim(B)}].$$
The thesis follows if we prove that 
\begin{equation}\label{goal}
\dim(B)p^{1-1/\dim(B)}<p/(100|B|)
\end{equation} 
as long as $c_1$ is chosen to be sufficiently small.

Set $h=dim(B)$, Equation \eqref{goal} is equivalent to
$$\log({h p^{1-1/h}})=\log(h)+\log{p}-1/h\log{p}<\log{p}-\log(100)-\log{|B|}$$
and so to
\begin{equation}\label{eq: r1}
    h \log(h) + h \log(100) + h\log|B| < \log p.
\end{equation}
We first consider the case where $R = c_1 (\log p)^{1/2}$. We note that due to Lemma \ref{lem:dimension-cardinality2} this relation is implied by
$$h\log{h}+2h\log{10}+h^2\log{5}<\log{p}
$$
which holds since we have assumed that $h < R = c_1 (\log p)^{1/2}$.

Let know consider the case where $h < R = c_1 \log p / \log |B|$. Here Equation \eqref{eq: r1} holds since for $c_1$ small enough we have
$$1/3 \log p \geq \max \left\{ h \log h, h \log (100), h \log |B| \right\}.$$
%We recall that, because of the definition of $R$, $h<c_1(\log{p})^{1/2}$ which implies that, if $c_1$ is sufficiently small,
%$$h\log{h}<c_1(\log{p})^{1/2}\log(c_1(\log{p})^{1/2})<\log{p}/3;$$
%$$2h\log{10}<2c_1(\log{p})^{1/2}\log(10)<\log{p}/3;$$
%$$h^2\log{5}<c_1^2(\log{p})\log{5}<\log{p}/3.$$
%Summing up these three relations we obtain Equation \eqref{goal} and %hence the thesis.
\end{proof}

\end{document}
